\SourceLeaf{357}%
Hunderttheile der Secunde geführt wird, und dass man sogar in der ganzen Zone von zwölf Graden, welche die Hülfstafel umfasst, die Berücksichtigung der Reductionen unterlassen kann, wenn man in der Rechnung nur Zehntel der Secunde notirt.

\pdfbookmark[0]{Anzeigen. Untersuchungen über Gegenstände der höhern Geodaesie. 1846.}{anzeigen-geodaesie-zweite}
\begin{center}
\rule{.85\textwidth}{.3pt}\par
Nachrichten von der K. Gesellschaft der Wissenschaften zu Göttingen.\quad 1846 September 28.\par
\rule{.85\textwidth}{.3pt}
\end{center}
\hspace*{\parindent}Am 1\textsuperscript{sten} September wurde von dem geh. Hofrath \textsc{Gauss} der Königlichen Gesellschaft der Wissenschaften eine Vorlesung überreicht mit der Überschrift:
\begin{center}\textit{Untersuchungen über Gegenstände der höhern Geodaesie, zweite Abhandlung},\end{center}
\noindent über deren Inhalt und Zusammenhang mit der ersten Abhandlung ein kurzer Bericht hier zu geben ist.

In der ersten Abhandlung war eine neue Methode, die geodaetischen Messungen zu behandeln, vorgetragen, deren Haupteigenthümlichkeit darin besteht, dass die meisten Berechnungen ganz oder fast ganz eben so geführt werden, als befände sich das Dreieckssystem nicht auf einer sphäroidischen, sondern auf einer Kugelfläche, und zwar ohne allen Abbruch für die äusserste Schärfe der Resultate. Eine der Hauptaufgaben im Gebiete der geodaetischen Rechnungen, nemlich aus der Grösse einer als geodaetische Linie auftretenden Dreiecksseite, der Breite des einen Endpunkts, und dem Azimuthe, unter welchem daselbst der andere Endpunkt erscheint, abzuleiten die Breite dieses andern Endpunkts, das dortige Azimuth der Dreiecksseite, und den Längenunterschied der beiden Punkte, reducirt sich bei jener Behandlungsweise auf die blosse Auflösung eines sphärischen Dreiecks. Ein Paar Seiten sind gleichwohl dieser Aufgabe in der erwähnten Abhandlung aus dem Grunde gewidmet, weil die gewöhnlichen Formeln der sphärischen Trigonometrie, wenn man nicht zu mehrziffrigen Logarithmen greifen will, nicht immer ausreichen würden, den Resultaten eine ganz genügende Schärfe zu geben, und deshalb gewisse Umformungen jener Formeln nothwendig werden. Ausserdem aber verstattet der Umstand, dass die Seiten solcher Dreiecke, deren Winkel wirklich gemessen werden, immer in Vergleich zu den Dimensionen des gan%
